% 图契约：从无线电波到伽马射线的频率升高与波长减小次序；几何、方向和公式与正文同一约定。
\documentclass[tikz,border=5pt]{standalone}
\usepackage{fontspec,xeCJK,amsmath,amssymb}
\IfFontExistsTF{Times New Roman}{\setmainfont{Times New Roman}}{\setmainfont{TeX Gyre Termes}}
\IfFontExistsTF{Heiti SC}{\setCJKmainfont[AutoFakeBold=1.5]{Heiti SC}}{\setCJKmainfont[AutoFakeBold=1.5]{FandolHei-Regular}}
\usetikzlibrary{arrows.meta,calc,patterns}
\definecolor{ink}{HTML}{243745}
\definecolor{blue}{HTML}{275DAB}
\definecolor{green}{HTML}{177C63}
\definecolor{red}{HTML}{BD4935}
\definecolor{orange}{HTML}{AA6C16}
\definecolor{purple}{HTML}{715393}
\begin{document}
\begin{tikzpicture}[x=1cm,y=1cm,font=\fontsize{9.5}{12}\selectfont,
v/.style={-{Latex[length=2mm,width=1.3mm]},line width=.85pt},
axis/.style={-{Latex[length=1.7mm,width=1.1mm]},ink!65,line width=.55pt},
aux/.style={ink!45,dashed,line width=.55pt},
lab/.style={fill=white,inner sep=1.8pt},
card/.style={draw=ink!25,fill=blue!3,rounded corners=2pt,text width=6.1cm,minimum height=.85cm,inner sep=6pt,align=center}]
\path[use as bounding box](-8,-3.6)rectangle(8,3.65);
\node[font=\bfseries\fontsize{10.5}{13}\selectfont,text width=15.1cm,align=center]at(0.00000,3.22000){M7-H28　从无线电波到伽马射线的频率升高与波长减小次序};
\begin{scope}[shift={(-4,0)}]
\node[]at(0.00000,2.30000){无线电};
\node[]at(0.00000,1.62000){微波};
\node[]at(0.00000,0.94000){红外};
\node[]at(0.00000,0.26000){可见光};
\node[]at(0.00000,-0.42000){紫外};
\node[]at(0.00000,-1.10000){X射线};
\node[]at(0.00000,-1.78000){$\gamma$射线};
\draw[v,red](-2.00000,2.50000)--(-2.00000,-2.10000);
\node[text=red]at(-2.10000,-2.50000){频率增大};
\draw[v,blue](2.00000,-2.10000)--(2.00000,2.50000);
\node[text=blue]at(2.00000,2.85000){波长增大};
\end{scope}
\node[card]at(4.00000,1.25000){$c=\lambda f$};
\node[card]at(4.00000,-1.05000){$\text{频率升高：无线电}\to\text{光}\to\text{X射线、}\gamma\text{射线}$};
\end{tikzpicture}\end{document}
